% Ficha HU-29. Criterios del backlog de la iteracion 5. El reporte no
% calcula nada nuevo: lee lo que dejo el punto de control.
\def\HUNro{29}
\def\HUSprint{5}
\def\HUTitulo{Reportes > Asistencia}
\def\HUModulo{M6}
\def\HUReq{12}
\def\HUPuntos{3}

\def\HUComo{docente}
\def\HUQuiero{obtener la asistencia de un examen}
\def\HUPara{documentar quiénes rindieron}

\def\HUDescripcion{%
Terminado el examen, el docente necesita el papel que dice quiénes rindieron. Hasta ahora
eso salía de una lista firmada a mano, con la letra de cada uno y los apellidos
incompletos.

El reporte lista a los estudiantes que registraron su ingreso, con la hora que puso el
servidor, el ambiente en el que entraron y el medio con que se los identificó. A quien
fue expulsado (HU-24) lo muestra igual ---estuvo en el aula--- pero marcado como tal, y
arrastra la observación que se anotó al ingresar (HU-26).

Se genera por examen y, cuando el examen se rinde en varias aulas, se lo puede acotar a
una. Los ingresos anulados (HU-21) no figuran: la anulación existe justamente para
deshacer un registro equivocado.

El reporte es de solo lectura, como todos los del módulo, y se exporta a PDF y a hoja de
cálculo (HU-34).}

\def\HUCriterios{%
\item Dado un examen con ingresos registrados, Cuando se pide el reporte, Entonces lista
      a los estudiantes que ingresaron con su hora.
\item Dado que el examen se rindió en varias aulas, Cuando se acota el reporte a una,
      Entonces solo figuran los estudiantes de esa aula.
\item Dado que un ingreso fue anulado, Cuando se pide el reporte, Entonces ese estudiante
      no figura.
\item Dado que un estudiante fue expulsado, Cuando se lo ve en el reporte, Entonces
      figura con su situación.
\item Dado que se anotó una observación al ingresar, Cuando se ve el reporte, Entonces la
      observación acompaña a esa fila.
\item Dado que el rol es el del docente, Cuando pide el reporte de su asignatura, Entonces
      el sistema se lo entrega; y le niega los consolidados.}

\def\HUPruebas{%
\item \texttt{ReportesTest::the\_attendance\_report\_lists\_who\_entered\_with\_their\_time}
      (criterios 1 y 5).
\item \texttt{ReportesTest::the\_attendance\_report\_can\_be\_narrowed\_to\_one\_room}
      (criterio 2).
\item \texttt{ReportesTest::the\_attendance\_report\_ignores\_annulled\_entries}
      (criterio 3).
\item \texttt{ReportesTest::the\_teacher\_reaches\_the\_reports\_of\_the\_subject\_only}
      (criterio 6).
\item \texttt{Reportes.test.jsx}: la pantalla muestra el reporte con sus filtros y su
      fecha de emisión (criterios 1 y 2).}

% =====================================================================
%  FICHA DE HISTORIA DE USUARIO  --  cuerpo comun del proyecto guia
%
%  Cada huNN.tex define antes de incluir este cuerpo:
%    \HUNro, \HUTitulo, \HUModulo, \HUReq, \HUPuntos, \HUSprint,
%    \HUComo, \HUQuiero, \HUPara, \HUDescripcion, \HUCriterios,
%    \HUPruebas  y, cuando la interfaz ya existe, los mockups en
%    mockups/huNN-escritorio.png y mockups/huNN-movil.png
%
%  Los criterios se redactan en formato Dado / Cuando / Entonces.
% =====================================================================
\documentclass[11pt,letterpaper]{article}
% Preambulo de las fichas del proyecto guia. Autocontenido: no depende del
% repositorio de documentacion, solo del logo que vive en logo/.
\usepackage[utf8]{inputenc}
\usepackage[T1]{fontenc}
\usepackage[spanish,es-noquoting]{babel}
\usepackage[margin=2.2cm,headheight=24pt,headsep=10pt]{geometry}
\usepackage{array}
\usepackage{longtable}
\usepackage{graphicx}
\usepackage{fancyhdr}
\usepackage{parskip}
\usepackage{enumitem}
\usepackage[breakable,skins]{tcolorbox}
\usepackage{needspace}

\newcolumntype{L}[1]{>{\raggedright\arraybackslash}p{#1}}
\newcolumntype{C}[1]{>{\centering\arraybackslash}p{#1}}
\setlength{\emergencystretch}{3em}

\newcommand{\NombreLargo}{NexaSoft S.R.L.}
\newcommand{\NombreProyecto}{Punku}
\newcommand{\Sistema}{Sistema informático para la gestión, validación y control del
ingreso de estudiantes durante la realización de exámenes masivos}

\newcommand{\LogoHorizontal}[1]{%
  \IfFileExists{logo/nexasoft_horizontal.png}{%
    \includegraphics[width=#1]{logo/nexasoft_horizontal}}{\textbf{[\,logo\,]}}}


\pagestyle{fancy}
\fancyhf{}
\fancyhead[L]{\raisebox{-3pt}{\LogoHorizontal{2.5cm}}}
\fancyhead[R]{\footnotesize Ficha de historia de usuario \,$\cdot$\, \emph{sprint} \HUSprint}
\fancyfoot[C]{\footnotesize \thepage}
\renewcommand{\headrulewidth}{0.4pt}

\setlength{\arrayrulewidth}{0.8pt}
\setlength{\tabcolsep}{4pt}

\newcommand{\seccion}[1]{%
  \vspace{0.6em}\needspace{6\baselineskip}\noindent{\Large\bfseries #1}\par
  \vspace{0.2em}}

% Caja que puede partirse entre paginas, para listas largas.
\newtcolorbox{CajaAncha}{breakable, enhanced jigsaw, sharp corners,
  colback=white, colframe=black, boxrule=0.8pt,
  width=\dimexpr16.6cm+2\tabcolsep+2\arrayrulewidth+5pt\relax,
  left=6pt, right=6pt, top=4pt, bottom=4pt}

% Mockup: muestra la captura si existe; si no, deja constancia de que la
% interfaz todavia no esta construida.
\newcommand{\CajaMockup}[2]{%
  \IfFileExists{mockups/#1.png}{%
    \includegraphics[width=#2]{mockups/#1}}{%
    \parbox[c][4.5cm][c]{#2}{\centering\itshape
      Pendiente: se genera al terminar la interfaz\\[2pt]
      \texttt{mockups/#1.png}}}}

\begin{document}

\noindent{\Large\bfseries Historia de Usuario \HUNro}
\vspace{0.3em}

\begin{center}
\begin{tabular}{|C{1.7cm}|L{13.0cm}|C{1.5cm}|}
\hline
\rule{0pt}{1.4em}\textbf{HU-\HUNro} &
\textbf{TÍTULO:} \HUTitulo &
{\footnotesize\HUModulo\ $\cdot$ Req. \HUReq\ $\cdot$ \HUPuntos\ pts} \\
\hline
\end{tabular}
\end{center}

\vspace{0.2em}
\begin{center}
\begin{tabular}{|L{16.6cm}|}
\hline
\rule{0pt}{1.2em}%
\begin{minipage}[t]{16.0cm}
\vspace{2pt}
\textbf{Como} \HUComo\par
\textbf{Quiero} \HUQuiero\par
\textbf{Para} \HUPara
\vspace{4pt}
\end{minipage} \\
\hline
\end{tabular}
\end{center}

\seccion{Descripción}

\begin{center}
\begin{CajaAncha}
\HUDescripcion
\end{CajaAncha}
\end{center}

\seccion{Mockups}

\begin{center}
\begin{tabular}{|C{11.0cm}|C{5.4cm}|}
\hline
\rule{0pt}{1.2em}{\footnotesize\bfseries Escritorio} & {\footnotesize\bfseries Móvil} \\
\CajaMockup{hu\HUNro-escritorio}{10.6cm} &
\CajaMockup{hu\HUNro-movil}{4.4cm} \\[4pt]
\hline
\end{tabular}
\end{center}

\ifx\HUMockupsExtra\undefined\else
\vspace{0.3em}
\begin{center}
\begin{tabular}{|C{8.2cm}|C{8.2cm}|}
\hline
\multicolumn{2}{|C{16.6cm}|}{%
  \rule{0pt}{1.2em}{\footnotesize\bfseries Otras pantallas de la historia}} \\
\HUMockupsExtra \\[4pt]
\hline
\end{tabular}
\end{center}
\fi
% Las demas pantallas en pareja: cada escritorio con su celular al lado. La
% tabla sigue en la pagina siguiente si no cabe.
\ifx\HUMockupsPares\undefined\else
\vspace{0.3em}
\begin{longtable}{|C{11.0cm}|C{5.4cm}|}
\hline
\multicolumn{2}{|C{16.6cm}|}{%
  \rule{0pt}{1.2em}{\footnotesize\bfseries Otras pantallas de la historia}} \\
\hline
\rule{0pt}{1.2em}{\footnotesize\bfseries Escritorio} & {\footnotesize\bfseries Móvil} \\
\hline
\endfirsthead
\hline
\rule{0pt}{1.2em}{\footnotesize\bfseries Escritorio} & {\footnotesize\bfseries Móvil} \\
\hline
\endhead
\HUMockupsPares
\end{longtable}
\fi

\seccion{Criterios de aceptación}

\begin{center}
\begin{CajaAncha}
\begin{enumerate}[leftmargin=1.6em, itemsep=3pt, topsep=0.2em]
\HUCriterios
\end{enumerate}
\end{CajaAncha}
\end{center}

\seccion{Pruebas automatizadas}

\begin{center}
\begin{CajaAncha}
% Los nombres de archivo en monoespaciada no se parten: se prefiere el
% margen derecho irregular antes que espacios enormes entre palabras.
\raggedright
\begin{itemize}[leftmargin=1.4em, itemsep=2pt, topsep=0.2em]
\HUPruebas
\end{itemize}
\end{CajaAncha}
\end{center}

\end{document}

